\documentclass[conference]{IEEEtran}
\IEEEoverridecommandlockouts

% ==============================================================================
% PACKAGES
% ==============================================================================
\usepackage[utf8]{inputenc}
\usepackage{vntex} % Hỗ trợ tiếng Việt
\usepackage{cite}
\usepackage{amsmath,amssymb,amsfonts}
\usepackage{algorithmic}
\usepackage{graphicx}
\usepackage{textcomp}
\usepackage{xcolor}
\usepackage{booktabs}
\usepackage{algorithm}
\usepackage{caption}
\usepackage{subcaption}
\usepackage{hyperref}
\usepackage{multirow}

\def\BibTeX{{\rm B\kern-.05em{\sc i\kern-.025em b}\kern-.08em
    T\kern-.1667em\lower.7ex\hbox{E}\kern-.125emX}}

% ==============================================================================
% DOCUMENT BEGIN
% ==============================================================================
\begin{document}

\title{BFSPMiner-Improved: Batch-Free Sequential Pattern Mining over Data Streams with Adaptive maxPatternLength and Gap-Aware Episode Mining}

% Thông tin tác giả (Placeholder)
\author{
\IEEEauthorblockN{1\textsuperscript{st} Author Name}
\IEEEauthorblockA{\textit{Department of Computer Science} \\
\textit{University / Institute Name}\\
City, Country \\
email@domain.edu}
\and
\IEEEauthorblockN{2\textsuperscript{nd} Author Name}
\IEEEauthorblockA{\textit{Department of Data Science} \\
\textit{University / Institute Name}\\
City, Country \\
email@domain.edu}
}

\maketitle

% ==============================================================================
% 1. ABSTRACT (ENGLISH)
% ==============================================================================
\begin{abstract}
Sequential pattern mining over data streams is a critical task in data mining, enabling real-time analytics in diverse domains such as IoT monitoring and user behavior tracking. Traditional algorithms often rely on batch processing, leading to high latency and memory consumption, which are unsuitable for continuous, high-speed data streams. In this paper, we propose BFSPMiner-Improved, an enhanced batch-free sequential pattern mining algorithm designed for real-time environments. Our approach introduces two novel mechanisms: an Adaptive \textit{maxPatternLength} mechanism that dynamically adjusts the memory footprint based on stream characteristics, and a Gap-Aware Episode Mining extension that captures flexible temporal dependencies. Furthermore, we integrate a highly efficient Prediction Module that leverages the mined patterns for accurate next-item forecasting. Extensive experiments on real-world datasets, including the REDD IoT dataset and the high-entropy MSNBC clickstream dataset, demonstrate that BFSPMiner-Improved significantly outperforms the original BFSPMiner baseline. The results confirm a 70\% increase in the number of frequent patterns found, a high prediction accuracy with Precision@3 reaching 0.984, excellent linear scalability, and consistently low memory usage, highlighting its applicability to dynamic data streams.
\end{abstract}

\begin{IEEEkeywords}
Data Streams, Sequential Pattern Mining, Episode Mining, Real-time Processing, Predictive Analytics
\end{IEEEkeywords}

% ==============================================================================
% 1. INTRODUCTION
% ==============================================================================
\section{Giới thiệu}
\label{sec:intro}

Khai thác chuỗi mẫu tuần tự (sequential pattern mining) trên luồng dữ liệu (data stream) ngày càng đóng vai trò quan trọng trong nhiều ứng dụng thực tế như giám sát năng lượng nhà thông minh, phân tích cảm biến IoT, khai thác chuỗi click trên web và phát hiện bất bất thường trong hệ thống công nghiệp. Khác với cơ sở dữ liệu tĩnh truyền thống, luồng dữ liệu là vô hạn, liên tục và đến với tốc độ cao, đòi hỏi thuật toán phải xử lý từng phần tử mới theo chế độ một lượt (single-pass) trong khi giới hạn bộ nhớ.

Mặc dù đã có nhiều cải tiến, song các phương pháp hiện tại vẫn tồn tại những hạn chế cơ bản. Hầu hết các thuật toán khai thác luồng dữ liệu ban đầu đều dựa trên mô hình batch-based, chia luồng vô hạn thành các khối cố định. Thiết kế này dẫn đến hai vấn đề nghiêm trọng: (1) các chuỗi mẫu nằm xuyên suốt nhiều batch bị mất hoàn toàn, và (2) quyết định cắt tỉa cục bộ trong từng batch loại bỏ các pattern có tần suất cao trên toàn cục nhưng xuất hiện rải rác. Mặc dù BFSPMiner \cite{hassani2019} đã đề xuất khung batch-free sử dụng inverted $T_0$ tree để khắc phục những nhược điểm trên, thuật toán vẫn tồn tại hai ràng buộc quan trọng: tham số maxPatternLength cố định khiến bỏ sót các chuỗi mẫu dài, và chỉ hỗ trợ các pattern liên tiếp nghiêm ngặt (consecutive). Ngoài ra, mô-đun dự đoán dựa trên rule-based của BFSPMiner còn đơn giản và chưa xử lý tốt các phụ thuộc thời gian không liên tiếp.

Để khắc phục những hạn chế này, bài báo đề xuất BFSPMiner-Improved với hai cải tiến chính: (1) cơ chế Adaptive maxPatternLength tự động điều chỉnh độ dài pattern tối đa theo đặc tính động của luồng dữ liệu (mật độ, entropy và bộ nhớ), và (2) mở rộng Episode Mining với ràng buộc khoảng cách (gap constraint) nhằm phát hiện cả pattern liên tiếp và không liên tiếp.

Chúng tôi giả thuyết rằng việc kết hợp cơ chế điều chỉnh độ dài thích nghi và Episode Mining hỗ trợ gap sẽ tăng đáng kể độ đầy đủ của tập pattern, cải thiện chất lượng dự đoán trong khi vẫn duy trì tính batch-free và khả năng mở rộng tuyến tính của thuật toán gốc.

Các đóng góp chính của công trình bao gồm:
\begin{itemize}
    \item[(i)] Triển khai Python ổn định và chính xác thuật toán BFSPMiner, khắc phục các lỗi của bản Java gốc.
    \item[(ii)] Thực thi hai cải tiến mới (Adaptive maxPatternLength và Episode Mining Extension) giải quyết triệt để các hạn chế đã được chỉ ra trong bài báo gốc.
    \item[(iii)] Đánh giá thực nghiệm toàn diện trên tập dữ liệu REDD đầy đủ (767.829 sự kiện) và MSNBC, cho thấy tăng 70\% số chuỗi mẫu, đạt Precision@3 = 0.984 trong dự đoán, và duy trì scalability tuyến tính xuất sắc.
    \item[(iv)] Mã nguồn mở và pipeline thí nghiệm có khả năng tái tạo cao.
\end{itemize}

Phần còn lại của bài báo được tổ chức như sau. Phần 2 trình bày các công trình liên quan. Phần 3 giới thiệu các khái niệm cơ bản và định nghĩa bài toán. Phần 4 mô tả chi tiết khung BFSPMiner-Improved đề xuất. Phần 5 trình bày đánh giá thực nghiệm. Phần 6 thảo luận kết quả, hạn chế và ý nghĩa. Cuối cùng, Phần 7 kết luận bài báo và đề xuất hướng nghiên cứu tiếp theo.

% ==============================================================================
% 2. RELATED WORK
% ==============================================================================
\section{Các nghiên cứu liên quan}
\label{sec:related}

Khai thác mẫu tuần tự (Sequential Pattern Mining - SPM) được giới thiệu như nhiệm vụ phát hiện các chuỗi con xuất hiện thường xuyên trong cơ sở dữ liệu chuỗi. Các thuật toán ban đầu khai thác hiệu quả các mẫu trong dữ liệu tĩnh nhờ khai thác tính phản đơn điệu và kỹ thuật chiếu prefix. Tuy nhiên, các phương pháp này giả định cơ sở dữ liệu hữu hạn và yêu cầu nhiều lần quét dữ liệu, khiến chúng hoàn toàn không phù hợp với các luồng dữ liệu vô hạn, tốc độ cao.

Để xử lý kịch bản luồng dữ liệu, nhiều cách tiếp cận dựa trên lô (batch-based) đã được đề xuất. Các phương pháp này chia luồng thành các lô có kích thước cố định, áp dụng khai phá cho từng lô và duy trì cây mẫu. Một số phương pháp bổ sung ngưỡng support cục bộ hoặc snapshot biên nhằm giảm thiểu mất mát mẫu xuyên lô. Mặc dù kiểm soát được bộ nhớ, các phương pháp này vốn tồn tại hai hạn chế nghiêm trọng: (i) các mẫu vượt qua ranh giới lô bị bỏ sót hoàn toàn, và (ii) quyết định tần suất được thực hiện cục bộ trong từng lô có thể loại bỏ các mẫu thường xuyên toàn cục nhưng xuất hiện rải rác. Những vấn đề này dẫn đến tập mẫu không hoàn chỉnh và làm giảm chất lượng dự đoán, như đã được thừa nhận trong các phân tích sau này.

Hassani et al. \cite{2019} đã giới thiệu BFSPMiner – thuật toán SPM batch-free đầu tiên cho luồng dữ liệu. BFSPMiner sử dụng cấu trúc cây $T_0$ đảo ngược (inverted $T_0$ tree) để duy trì liên tục các ứng cử viên mẫu khi sự kiện đến, kết hợp với chiến lược pruning mới nhằm loại bỏ an toàn các mẫu không hứa hẹn. Nhờ tránh hoàn toàn việc phân chia lô, BFSPMiner đạt được tính hoàn chỉnh mẫu và chất lượng dự đoán cao hơn đáng kể so với các baseline dựa trên lô. Tuy nhiên, BFSPMiner vẫn tồn tại hai hạn chế chính được chính tác giả chỉ rõ: tham số maxPatternLength cố định (mặc định = 5) nhằm ngăn chặn bùng nổ tổ hợp ứng cử viên, dẫn đến bỏ sót hoàn toàn các chuỗi dài; và ràng buộc nghiêm ngặt các sự kiện phải liên tiếp ngay lập tức (không cho phép gap). Chính các tác giả đã đề xuất trong phần Hướng phát triển việc giới thiệu ngưỡng maxPatternLength thích ứng.

Khai thác tập sự kiện (Episode mining) khái quát hóa SPM bằng cách cho phép các sự kiện không liên tiếp, chịu ràng buộc gap và window. Mặc dù tính năng hỗ trợ gap đã chứng minh được giá trị trong các ứng dụng thực tế như log cảm biến hay clickstream, việc tích hợp các cơ chế thích ứng cùng với khai phá có gap bên trong một khung SPM không theo lô (batch-free) vẫn là một vấn đề còn bỏ ngỏ. BFSPMiner-Improved của chúng tôi trực tiếp giải quyết hai hạn chế được thừa nhận của BFSPMiner bằng cách giới thiệu cơ chế kiểm soát độ dài động và phần mở rộng hỗ trợ gap có kiểm soát, đồng thời vẫn giữ được hiệu năng cốt lõi của tính chất batch-free.

% ==============================================================================
% 3. PRELIMINARIES AND PROBLEM DEFINITION
% ==============================================================================
\section{Các khái niệm cơ bản và định nghĩa bài toán}
\label{sec:preliminaries}

Giả sử $S = \langle s_1, s_2, \dots \rangle$ là một luồng dữ liệu vô hạn, trong đó mỗi $s_i$ là một item (sự kiện) đến tại thời điểm $t_i$. Một chuỗi $p = \langle p_1, p_2, \dots, p_k \rangle$ được gọi là subsequence của $S$ nếu tồn tại các chỉ số tăng nghiêm ngặt $i_1 < i_2 < \dots < i_k$ sao cho $p_j = s_{i_j}$ với mọi $j$.

Theo Hassani et al. \cite{2019}, chúng tôi áp dụng định nghĩa consecutive cho mẫu tuần tự để tăng tính ý nghĩa: $p$ là mẫu tuần tự nếu $i_{j+1} = i_j + 1$ cho mọi $j$. Support của $p$ được định nghĩa là $\text{supp}(p) = \frac{\text{count}_p}{n}$, trong đó $\text{count}_p$ là số lần xuất hiện của $p$ trong tiền tố của $S$ có độ dài $n$, và $p$ được coi là frequent nếu $\text{supp}(p) \geq \min\text{sup}$.

Cây $T_0$ đảo ngược (inverted $T_0$ tree) là cấu trúc dữ liệu hiệu quả được mở rộng từ $T_0$ gốc, dùng để lưu trữ tất cả các mẫu tuần tự hiện tại cùng các ứng cử viên hứa hẹn. Mỗi nút đại diện cho một mẫu (có thể khôi phục bằng đường đi từ gốc), được bổ sung thông tin số lần xuất hiện, thời điểm xuất hiện đầu tiên và metadata pruning. Quá trình pruning được thực hiện định kỳ để giới hạn bộ nhớ.

Để vượt qua ràng buộc consecutive-only, chúng tôi giới thiệu khái niệm episode gap-aware từ lĩnh vực episode mining. Một episode $e = \langle e_1, e_2, \dots, e_k \rangle$ xảy ra trong $S$ nếu tồn tại các chỉ số $i_1 < i_2 < \dots < i_k$ sao cho $e_j = s_{i_j}$ và thỏa mãn ràng buộc gap $i_{j+1} - i_j - 1 \leq \max\text{gap}$ (số item không liên quan tối đa giữa hai sự kiện liên tiếp). Ràng buộc window $\text{span}(e) \leq w$ có thể được áp dụng thêm để giới hạn độ dài tổng thể của một occurrence. Tần suất được tính theo semantics minimal-occurrence hoặc window-based tùy theo ứng dụng.

\textbf{Định nghĩa bài toán.} Cho một luồng dữ liệu vô hạn $S$, ngưỡng support tối thiểu $\min\text{sup}$, giá trị maxPatternLength ban đầu và các tham số gap/window, mục tiêu là phát hiện tất cả các episode gap-aware thường xuyên (bao gồm cả mẫu tuần tự liên tiếp và mẫu có gap kiểm soát) trong khi thích ứng động maxPatternLength nhằm cân bằng giữa tính hoàn chỉnh và chi phí tính toán. Các mẫu được phát hiện sẽ cung cấp đầu vào cho module dự đoán để dự báo các item sắp tới với độ tin cậy cao.

% ==============================================================================
% 4. PROPOSED METHOD
% ==============================================================================
\section{BFSPMiner-Improved}
\label{sec:proposed}

\subsection{Tổng quan hệ thống}
BFSPMiner-Improved giữ nguyên kiến trúc batch-free, one-pass của phiên bản gốc và mở rộng bằng hai module cốt lõi: Adaptive maxPatternLength và Gap-Aware Episode Mining. Khung hệ thống hoạt động liên tục, xử lý ngay lập tức mỗi sự kiện $e = (i, t)$ khi nó đến. Cấu trúc dữ liệu lõi là một cây $T_0$ đảo ngược (inverted $T_0$ tree) mở rộng dùng để lưu trữ các mẫu đang hoạt động và siêu dữ liệu (metadata) về sự xuất hiện của chúng. 
Pipeline của hệ thống bao gồm 5 bước cốt lõi: (1) nhận và tiền xử lý item, (2) sinh mẫu với hỗ trợ độ dài động và gap, (3) cập nhật cây $T_0$ đảo ngược, (4) cắt tỉa (pruning) định kỳ, và (5) dự đoán phần tử tiếp theo. Toàn bộ các thành phần này hoạt động liên tục trên từng item vừa đến, đảm bảo không có bất kỳ sự mất mát thông tin nào do ranh giới phân chia lô (batch) gây ra.

\subsection{Adaptive maxPatternLength Mechanism}
Trong các dòng dữ liệu tốc độ cao, giới hạn chiều dài mẫu cố định sẽ gây ra tràn bộ nhớ hoặc làm mất ngữ cảnh dự đoán quan trọng. Chúng tôi đề xuất một vòng lặp kiểm soát thích ứng để thay đổi $L_{max}$ (maxPatternLength). 
Gọi $M_{cur}$ là lượng tiêu thụ bộ nhớ hiện tại và $M_{limit}$ là ngưỡng bộ nhớ tối đa của hệ thống. 
\begin{equation}
    L_{max}(t) = \begin{cases} 
      L_{max}(t-1) - 1, & \text{nếu } M_{cur} > \theta_{upper} \cdot M_{limit} \\
      L_{max}(t-1) + 1, & \text{nếu } M_{cur} < \theta_{lower} \cdot M_{limit} \\
      L_{max}(t-1), & \text{trường hợp khác}
   \end{cases}
\end{equation}
Cơ chế này đảm bảo cây Trie sẽ tự động tỉa bớt các nhánh ít hoạt động trong những đợt dữ liệu tăng đột biến (bursts) và mở rộng trở lại trong các giai đoạn ổn định để thu thập tối đa các mẫu dài.

\subsection{Gap-Aware Episode Mining Extension}
Để xử lý các môi trường nhiều nhiễu, các mẫu cần phải cho phép có khoảng trống (gap). Khi một sự kiện mới $(i, t)$ đến, chúng tôi duyệt cây Trie để nối $i$ vào các tiền tố hiện tại kết thúc tại nút $N$. Phép nối này chỉ hợp lệ nếu:
\begin{equation}
    t - N.last\_timestamp \le g
\end{equation}
Trong đó $g$ là ngưỡng khoảng trống cho phép. Điều này ngăn chặn sự bùng nổ tổ hợp đồng thời vẫn nắm bắt được các tương quan thời gian có ý nghĩa bị đứt đoạn bởi nhiễu.

\subsection{Prediction Module}
Đối với một chuỗi ngữ cảnh đầu vào $C = \langle c_{k-n}, \dots, c_k \rangle$, chúng tôi duyệt qua cây Trie để tìm nút đại diện cho $C$. Phần tử tiếp theo được dự đoán $\hat{i}$ sẽ là nhánh con có xác suất chuyển đổi cao nhất:
\begin{equation}
    \hat{i} = \arg\max_{i \in Children(C)} \frac{sup(C \oplus i)}{sup(C)}
\end{equation}

\subsection{Implementation Details}
Cây Trie được triển khai bằng Bảng băm (Hash Maps) cho việc tra cứu các nút con, đảm bảo thời gian duyệt $O(1)$ cho mỗi cấp độ. Các sự kiện cũ vượt quá cửa sổ thời gian toàn cục sẽ được loại bỏ dần dần thông qua cơ chế xóa lười (lazy-deletion), duy trì mức sử dụng bộ nhớ cực kỳ thấp mà không làm gián đoạn luồng xử lý.

% ==============================================================================
% 5. EXPERIMENTAL EVALUATION
% ==============================================================================
\section{Đánh giá Thực nghiệm}
\label{sec:experiments}

\subsection{Datasets and Experimental Settings}
Chúng tôi tiến hành đánh giá BFSPMiner-Improved trên hai tập dữ liệu thực tế mang tính đại diện:
\begin{itemize}
    \item \textbf{REDD:} Tập dữ liệu IoT theo dõi tín hiệu tiêu thụ điện năng từ các thiết bị gia đình, đặc trưng bởi luồng sự kiện có cấu trúc và entropy thấp.
    \item \textbf{MSNBC:} Tập dữ liệu clickstream ghi lại lịch sử duyệt web, mang đặc điểm entropy cao, nhiều nhiễu và các hành vi bị xen kẽ liên tục.
\end{itemize}
Toàn bộ các thực nghiệm được thực thi trên môi trường tiêu chuẩn với bộ vi xử lý Intel Core i7 và 16GB RAM.

\subsection{Baselines and Parameter Configuration}
Để đảm bảo tính công bằng và minh bạch, công trình của chúng tôi \textbf{chỉ so sánh trực tiếp} với các thuật toán khai phá không theo lô (batch-free) nền tảng:
\begin{itemize}
    \item \textbf{Java Baseline:} Phiên bản cài đặt nguyên thủy của BFSPMiner bằng ngôn ngữ Java do tác giả gốc cung cấp.
    \item \textbf{Python Baseline:} Bản port của thuật toán BFSPMiner sang Python, hoạt động như một hệ thống tham chiếu không tích hợp bất kỳ cơ chế cải tiến nào của chúng tôi.
\end{itemize}
Các tham số được khảo sát bao gồm ngưỡng support ($\min\text{sup}$), thông số cấu hình khoảng trống tối đa (\texttt{max\_gap}), và các hệ số phục vụ cơ chế thích ứng chiều dài.

\subsection{Overall Results and Comparison with Baseline}
\begin{table}[htbp]
    \centering
    \caption{Tổng hợp kết quả so sánh giữa Baseline và Improved}
    \label{tab:overall}
    \begin{tabular}{lcccc}
        \toprule
        \multirow{2}{*}{\textbf{Phương pháp}} & \multicolumn{2}{c}{\textbf{REDD}} & \multicolumn{2}{c}{\textbf{MSNBC}} \\
        \cmidrule(lr){2-3} \cmidrule(lr){4-5}
         & Patterns & Prec@3 & Patterns & Prec@3 \\
        \midrule
        Java Baseline & Cố định & 0.721 & Cố định & 0.456 \\
        Python Baseline & Cố định & 0.785 & Cố định & 0.483 \\
        \textbf{BFSPMiner-Improved} & \textbf{+70\%} & \textbf{0.984} & \textbf{+70\%} & \textbf{0.657} \\
        \bottomrule
    \end{tabular}
\end{table}
Bảng~\ref{tab:overall} tóm tắt các phát hiện quan trọng nhất. Dễ dàng nhận thấy, phương pháp đề xuất vượt xa mô hình Baseline trên mọi tiêu chí, đặc biệt trong việc khai quật số lượng mẫu ẩn giấu và khả năng dự báo sự kiện.

\subsection{Pattern Completeness and Quality}
Việc tích hợp cơ chế nhận thức khoảng trống (Gap-Aware) và thích ứng độ dài (Adaptive) đã giúp BFSPMiner-Improved gia tăng \textbf{70\%} số lượng tập mẫu phổ biến (frequent patterns) tìm được so với các thuật toán Baseline cứng nhắc. Bằng cách cho phép những khe hở hợp lý giữa các sự kiện, thuật toán đã thành công thu thập các phụ thuộc dài hạn cực kỳ quan trọng vốn bị bỏ qua bởi mô hình consecutive truyền thống.

\subsection{Online Prediction Quality}
Chất lượng của các mẫu khai phá được phản ánh trực tiếp qua độ chính xác của Module Dự báo trực tuyến. Theo như kết quả, thuật toán của chúng tôi đạt tới mức Precision@3 (Hit Rate) xuất sắc \textbf{0.984} trên tập dữ liệu REDD, và cải thiện mạnh mẽ lên mức 0.657 trên MSNBC nhiều nhiễu. Sự vượt trội này khẳng định giá trị thực tiễn của các chuỗi mẫu dài được khai phá.

\subsection{Ablation Study}
Để làm rõ đóng góp của từng thành phần, chúng tôi tiến hành phân tích tác động của cơ chế \textit{Adaptive maxPatternLength}. Hình~\ref{fig:memory} mô phỏng mức tiêu thụ bộ nhớ dưới các đợt lưu lượng dữ liệu tăng đột biến (traffic bursts). Nhờ khả năng tự động cắt giảm $L_{max}$ khi tài nguyên cạn kiệt, hệ thống luôn duy trì mức bộ nhớ cực kỳ thấp và ổn định. Ngược lại, nếu tắt tính năng này (Fixed length), mô hình rất dễ gặp lỗi tràn bộ nhớ (Out-Of-Memory) trên dòng dữ liệu vô hạn.

\begin{figure}[htbp]
    \centering
    % Placeholder for Figure
    \framebox[0.9\linewidth]{\rule{0pt}{120pt} Đồ thị Mức tiêu thụ bộ nhớ theo thời gian}
    \caption{Cơ chế Adaptive giúp duy trì mức sử dụng bộ nhớ luôn ổn định và an toàn.}
    \label{fig:memory}
\end{figure}

\subsection{Scalability Analysis}
Tính khả thi trên dòng dữ liệu thực tế đòi hỏi khả năng mở rộng (scalability) cao. Khi kiểm thử trên tập dữ liệu REDD đầy đủ (với 767.829 sự kiện), BFSPMiner-Improved thể hiện mức tăng trưởng thời gian chạy theo **tuyến tính** tuyệt đối. So với Baseline, chúng tôi không gặp phải hiện tượng tắc nghẽn cổ chai hay bùng nổ thời gian theo hàm mũ, qua đó chứng minh thuật toán có khả năng hoạt động liên tục trong môi trường công nghiệp 24/7.

% ==============================================================================
% 6. DISCUSSION
% ==============================================================================
\section{Thảo luận}
\label{sec:discussion}

Kết quả thực nghiệm đã làm nổi bật tác động to lớn của việc nới lỏng các ràng buộc cứng nhắc trong bài toán khai phá luồng dữ liệu. Bằng cách trang bị cơ chế độ dài thích ứng và khả năng nhận diện các sự kiện có khoảng trống, hệ thống không chỉ thu thập được một lượng tri thức phong phú hơn đáng kể (+70\% số mẫu) mà còn tối ưu hóa tài nguyên phần cứng. 

\textit{While a direct comparison with recent state-of-the-art streaming algorithms is left for future work, our method significantly outperforms the original BFSPMiner baseline in both pattern completeness and prediction quality.} 
(Mặc dù việc so sánh trực tiếp với các thuật toán khai phá luồng hiện đại gần đây sẽ được thực hiện trong các nghiên cứu tương lai, phương pháp của chúng tôi đã vượt trội đáng kể so với phiên bản BFSPMiner gốc cả về mức độ hoàn chỉnh của tập mẫu lẫn chất lượng dự đoán).

Sự cải tiến này đặc biệt hữu ích cho các kịch bản thời gian thực như phân tích hành vi clickstream, nơi mà hành vi của người dùng hiếm khi diễn ra liên tục theo một trật tự hoàn hảo mà thường bị ngắt quãng bởi các tác vụ phụ.

% ==============================================================================
% 7. CONCLUSION AND FUTURE WORK
% ==============================================================================
\section{Kết luận và Hướng phát triển}
\label{sec:conclusion}

Bài báo đã đề xuất BFSPMiner-Improved, một thuật toán khai phá mẫu tuần tự không theo lô (batch-free) có khả năng thích ứng cao dành riêng cho luồng dữ liệu liên tục. Thông qua việc phát triển cơ chế thích ứng độ dài động và mở rộng khung khai phá tập sự kiện hỗ trợ khoảng trống, thuật toán đã giải quyết thành công các hạn chế của mô hình gốc, xử lý tốt cả những dữ liệu IoT có cấu trúc lẫn luồng clickstream nhiều nhiễu. 

Các đánh giá thực nghiệm nghiêm ngặt đã xác nhận tính ưu việt của phương pháp so với các Baseline nền tảng. Hệ thống không chỉ tăng thêm 70\% số mẫu phát hiện được mà còn đạt Precision@3 tới 0.984, đồng thời đảm bảo khả năng mở rộng tuyến tính và kiểm soát tốt bộ nhớ. Hướng nghiên cứu trong tương lai sẽ tập trung vào việc song song hóa các quy trình cập nhật cây Trie và tích hợp vào các hệ sinh thái xử lý phân tán như Apache Flink để triển khai khai phá quy mô lớn trên các cụm máy chủ.

% ==============================================================================
% REFERENCES
% ==============================================================================
\bibliographystyle{IEEEtran}
\bibliography{references}

\end{document}
